O controle inteligente de sistemas complexos, como ambientes fechados, pode ser realizado de forma eficiente utilizando lógica fuzzy~\cite{zadeh1965fuzzy,ross2010fuzzy}. Este relatório apresenta o desenvolvimento de dois sistemas fuzzy distintos: um sistema Mamdani para controle de ventilação em ambientes fechados~\cite{mamdani1975experiment,ross2010fuzzy} e um sistema Takagi-Sugeno para aproximação de funções não lineares~\cite{takagi1985fuzzy,sugeno1985industrial}. Ambos os sistemas foram implementados em Python, sem o uso de bibliotecas específicas de lógica fuzzy, visando promover o entendimento dos conceitos fundamentais e a flexibilidade na modelagem.

O sistema Mamdani tem como objetivo auxiliar na tomada de decisão sobre a intensidade da ventilação, promovendo conforto térmico e eficiência energética, a partir de variáveis ambientais e ocupacionais~\cite{ross2010fuzzy}. Já o sistema Takagi-Sugeno é utilizado para aproximar a função não linear $f(x) = e^{-x/5} \cdot \sin(3x) + 0.5 \cdot \sin(x)$ no intervalo $x \in [0, 10]$, explorando diferentes funções de pertinência, operadores fuzzy e técnicas de otimização dos consequentes~\cite{takagi1985fuzzy}.

O relatório está organizado da seguinte forma: na Seção de Metodologia são detalhados os modelos fuzzy, as funções de pertinência, as bases de regras e os métodos de inferência e defuzzificação utilizados. Em seguida, são apresentados os resultados das simulações, análises comparativas e discussões sobre o desempenho dos sistemas. Por fim, são apresentadas as conclusões e sugestões para trabalhos futuros.

